\chapter*{About These Notes}
\addcontentsline{toc}{chapter}{About These Notes}

\paragraph{Goal.} These notes accompany a four-hour course for readers who
know real analysis and mathematical statistics but have not met queueing
theory. The goal is narrow: to be able to prove, unaided, the four results
our research uses to model and analyse agentic LLM serving with queues.
Each lecture takes about an hour, and every concept is defined before it is
used. Measure-theoretic probability is not needed; we use random variables,
expectation, conditional expectation, the strong law of large numbers
(SLLN) and power series freely.

\paragraph{Background: why LLM serving is a queue.} An LLM request is
served in two phases. \emph{Prefill} reads the whole prompt in one pass and
produces the first token; the time from the request's arrival to its first
token is the \emph{time to first token} (TTFT). \emph{Decode} then produces
one token per step. Both phases keep the attention state of every prompt
token, the \emph{KV cache}, in device memory. When a later request starts
with the same tokens, the stored state can be reused (a \emph{hit}); if it
has been evicted from memory, it must be recomputed (a \emph{miss}). An
agent program is a \emph{session} that loops through
prefill $\to$ decode $\to$ tool call $\to$ prefill $\to\cdots$, and each
request (a \emph{turn}) is the previous prompt with a tool result appended.
Requests arrive at random, wait for a server (the accelerator), receive
service and leave. This is exactly the object of queueing theory.

\paragraph{Map.} The four results of our research, and what these notes
provide to prove them:

\begin{center}\small
\begin{tabular}{@{}p{0.3\linewidth}p{0.64\linewidth}@{}}
\toprule
Research result & Needed from these notes \\
\midrule
The decode stage prices only work &
  Little's law (\cref{thm:little}), birth--death processes
  (\cref{thm:birthdeath}), processor sharing and insensitivity
  (\cref{thm:insens}), \cref{thm:L-mono}, \cref{prop:decode-price} \\
The price of a miss &
  PASTA (\cref{thm:pasta}), residual service time (\cref{lem:residual}),
  the Pollaczek--Khinchine formula (\cref{thm:pk}), \cref{thm:price} \\
Finite population &
  closed systems, the arrival theorem (\cref{thm:arrival}), mean value
  analysis (\cref{thm:mva}), \cref{prop:finite} \\
Price-blind eviction keys &
  the covering knapsack, the threshold rule (\cref{thm:threshold}),
  \cref{prop:blind}; Campbell's theorem (\cref{thm:campbell}) for the
  memory estimate \\
\bottomrule
\end{tabular}
\end{center}

\paragraph{Plan.} \Cref{sec:L1} covers arrival processes and Little's law;
\cref{sec:L2} Markovian queues and processor sharing; \cref{sec:L3} the
M/G/1 queue, the Pollaczek--Khinchine formula and the price of a miss;
\cref{sec:L4} closed systems and memory; and \cref{sec:L5} assembles these tools into the theory of the paper, step by step. The exercises at the end of each
lecture are meant to be solved before the next one. Good introductory
textbooks are \cite{kleinrock1975,harchol2013}; \cite{asmussen2003} is the
rigorous reference.
